% Generated draft: requires mathematical and source review.
\documentclass[11pt]{article}
\usepackage[margin=25mm]{geometry}
\usepackage{fontspec}
\setmainfont{Latin Modern Roman}
\usepackage{amsmath,amssymb,mathtools,unicode-math}
\setmathfont{Latin Modern Math}
\usepackage{enumitem,needspace,xurl,hyperref}
\setlength{\parindent}{0pt}
\setlength{\parskip}{5pt}
\setlength{\emergencystretch}{3em}
\newcommand{\N}{\mathbb{N}}
\newcommand{\Z}{\mathbb{Z}}
\newcommand{\Q}{\mathbb{Q}}
\newcommand{\R}{\mathbb{R}}
\newcommand{\C}{\mathbb{C}}
\newcommand{\F}{\mathbb{F}}
\newcommand{\A}{\mathbb{A}}
\newcommand{\PP}{\mathbb{P}}
\newcommand{\E}{\mathbb{E}}
\newcommand{\eps}{\varepsilon}
\newcommand{\dd}{\,\mathrm{d}}
\begin{document}
\section*{UnsolvedMath ID: 20000732}
\textbf{Source problem code:} catalogue entry
\subsection{Definitions and mathematical statement}
The source's symbols $A$ and $B$ are not defined in the isolated problem text. The supplied research aid interprets them through Nahm sums. In the standard convention, for rational symmetric positive-definite $A\in M_r(\Q)$, $B\in\Q^r$, and $C\in\Q$, put
\[
 F_{A,B,C}(q)=\sum_{n\in\N_0^r}
 \frac{q^{\frac12n^{\mathsf T}An+B^{\mathsf T}n+C}}{(q)_{n_1}\cdots(q)_{n_r}},\qquad
 (q)_m=\prod_{j=1}^m(1-q^j),\quad(q)_0=1,
\]
with $q=e^{2\pi i\tau}$, $\operatorname{Im}\tau>0$. A scalar weight-zero modular function for a finite-index group $\Gamma$ is invariant under $\Gamma$, holomorphic on the upper half-plane here, and meromorphic at its cusps; the reported span bound additionally requires trivial multiplier.

\textbf{Problem 2.1.} For fixed $A$, bound the number of possible $B$, and bound the dimension of the vector space determined by the $B$. The source allows ineffective as well as effective bounds. Under the aid's interpretation, ``possible'' means admitting a modular Nahm sum for some $C$, and the vector space is the complex span of the resulting functions. The precise original surrounding definition is not supplied.

\textbf{Reported algebraic certificate.} Scalar modularity forces necessary Vlasenko--Zwegers equations
\[
 P_p^A(B)=0,\qquad\deg_B P_p^A\leq2p,
\]
and uniquely determines $C$. Their explicit coefficients, index convention and the formula for $C$ are absent from the summary. Over a field $K$ containing their algebraic coefficients, consider a finite initial ideal $I_m$ generated by the first $m$ equations. If its complex obstruction locus is zero-dimensional, its finite quotient length
\[
 D_m=\dim_K K[B_1,\ldots,B_r]/I_m
\]
bounds the number of distinct modular $B$ by $D_m$, with multiplicities included in the length. The empty-locus case has bound zero. The aid invokes Noetherianity to say that if the full necessary-equation locus is finite, some finite truncation supplies such a certificate.

This is conditional on zero-dimensionality of the obstruction locus. It does not prove that locus finite for every $A$, that every solution is modular, or that mere finiteness of the actually modular $B$ forces the larger obstruction locus to be finite. Effective use also needs the equations in exact computable form.

\textbf{Reported common-group span bound.} Fix one common finite-index group $\Gamma$ on which all functions under consideration are scalar weight-zero modular functions with trivial multiplier. For each cusp $\mathfrak a$, let $h_{\mathfrak a}$ be its width. The aid cites a Garoufalidis--Zagier valuation bound denoted by $C_0(A)$; its definition is not included. In the convention where
\[
 \operatorname{ord}_{\mathfrak a}F\geq h_{\mathfrak a}C_0(A),
\]
the corresponding effective pole divisor has coefficients $\max\{0,\lfloor-h_{\mathfrak a}C_0(A)\rfloor\}$. Riemann--Roch then yields the reported bound
\[
 \dim_\C\operatorname{span}\{F_{A,B,C}\}
 \leq1+\sum_{\mathfrak a\ \mathrm{cusp\ of}\ \Gamma}
 \max\{0,\lfloor-h_{\mathfrak a}C_0(A)\rfloor\}.
\]
The exact valuation theorem and $C_0(A)$ must be verified before numerical use. The result depends on the common group and its cusp widths as well as $A$; it is not an $A$-only bound when modular groups and multipliers vary freely. The original fixed-$A$ problem remains broader than these certificate and common-group reductions.
\subsection{Short English statement}
For the report's Nahm-sum interpretation, modular parameters satisfy necessary polynomial equations. A zero-dimensional finite subsystem gives an effective bound on their number, but its existence is conditional. Separately, a common scalar modular group and a cusp valuation bound control the dimension of the span; allowing arbitrary groups is a broader problem.
\subsection{Sources}
\begin{itemize}
\item Local UnsolvedMath registry, permanent ID $20000732$, \texttt{catalogue entry}; canonical \texttt{aim-analytic-number-theory-notes.json}, \texttt{notes[95]}.
\item AIM workshop ``Mock modular forms in combinatorics and arithmetic geometry,'' PDF Problem 2.1: \url{https://aimath.org/WWN/mockmodular/mockmodular.pdf}. The isolated extract omits the surrounding definitions of $A,B$; no PDF inspection is asserted here.
\item Vlasenko--Zwegers polynomial obstructions and Garoufalidis--Zagier cusp valuation bound, as named in the supplied summary; no complete theorem statements or bibliographic references are supplied.
\item Research attempt 1: \texttt{reduction}, run status \texttt{open}, internal rigor \texttt{proved}, low novelty confidence; current registry partially solved. Explicitly a machine-generated research aid, not peer reviewed; independently verify before citation. Original literature field empty.
\end{itemize}
\end{document}
